
\begin{figure}[htbp]
    \centering
    \begin{subfigure}[b]{0.4\textwidth}
        \centering
        \begin{tikzpicture}[scale=1.5]
            \tikzstyle{vertex}=[circle,fill=black!25,minimum size=17pt,inner sep=0pt]
            \tikzstyle{edge} = [draw,thick,-]
    
            \matrix[matrix of math nodes,
              %left delimiter={[},
              %right delimiter={]},
              %below = 1cm of A,
              nodes={minimum width=3em, minimum height=3em, outer sep=0pt},
              column sep=-\pgflinewidth,
              row sep=-\pgflinewidth] (m) { A, & B, & C \\
                                            B, & C, &   \\
                                            C, &    &   \\
                                            D, & C, &   \\
                                            E, & A, &   \\
                                        };
    
        \end{tikzpicture}
        \caption{Adjacency List representation}
        \label{figure:graph-adj-list}
    \end{subfigure}
    \hfill
    \begin{subfigure}[b]{0.4\textwidth}
        \centering
        \begin{tikzpicture}[>=stealth, shorten >=1pt, auto, node distance=3cm, thick, main node/.style={circle, draw, font=\sffamily\Large\bfseries}]

        % Render the adjacency matrix
        \matrix[matrix of math nodes,
                left delimiter={[},
                right delimiter={]},
                %below = 1cm of A,
                nodes={minimum width=3em, minimum height=3em, outer sep=0pt},
                column sep=-\pgflinewidth,
                row sep=-\pgflinewidth] (m) {   & A & B & C & D & E \\
                                                A & 0 & 1 & 1 & 0 & 0 \\
                                                B & 0 & 0 & 1 & 0 & 0 \\
                                                C & 0 & 0 & 0 & 0 & 0 \\
                                                D & 0 & 0 & 1 & 0 & 0 \\
                                                E & 1 & 0 & 0 & 0 & 0 \\
                                            };

        \end{tikzpicture}
        \caption{Adjacency Matrix representation: a $|V| \times |V|$ matrix where a value of $1$ indicates an edge and $0$ is no edge.}
        \label{figure:adjacency_matrix}
    \end{subfigure}
    
    \medskip
    
    \begin{subfigure}[b]{0.4\textwidth}
        \centering
        \begin{tikzpicture}[>=stealth, shorten >=1pt, auto, node distance=3cm, thick, main node/.style={circle, draw, font=\sffamily\Large\bfseries}]
        
        % Render the adjacency matrix
        \matrix[matrix of math nodes,
                    left delimiter={[},
                    right delimiter={]},
                    %below = 1cm of A,
                    nodes={minimum width=3em, minimum height=3em, outer sep=0pt},
                    column sep=-\pgflinewidth,
                    row sep=-\pgflinewidth] (m) {   & e_0 & e_1 & e_2 & e_3 & e_4\\ 
                                                A   &  -1 &   0 &   0 &   1 &  -1\\
                                                B   &   1 &  -1 &   0 &   0 &   0\\
                                                C   &   0 &   1 &   1 &   0 &   1\\
                                                D   &   0 &   0 &  -1 &   0 &   0\\
                                                E   &   0 &   0 &   0 &  -1 &   1\\
                                                };
        
        \end{tikzpicture}
        \caption{Incidence Matrix representation: an $|E| \times |V|$ matrix where a -1 is the source node and 1 is the destination of a given edge.}
        \label{figure:incidence_matrix}
    \end{subfigure}
    \hfill
    \begin{subfigure}[b]{0.4\textwidth}
        \centering
        \begin{tikzpicture}[>=stealth, shorten >=1pt, auto, node distance=3cm, thick, main node/.style={circle, draw, font=\sffamily\Large\bfseries}]
        
        % Render the adjacency matrix
            \matrix[matrix of math nodes,
                left delimiter={[},
                right delimiter={]},
                %below = 1cm of A,
                nodes={minimum width=3em, minimum height=3em, outer sep=0pt},
                column sep=-\pgflinewidth,
                row sep=-\pgflinewidth] (M1) { 
                                                1 & 1 & 1 & 1 & 1 &\\
                                                };
            \node[left] at (M1.west) {Value List..};                                    
        \end{tikzpicture}

        \begin{tikzpicture}[>=stealth, shorten >=1pt, auto, node distance=3cm, thick, main node/.style={circle, draw, font=\sffamily\Large\bfseries}]

            \matrix[matrix of math nodes,
                left delimiter={[},
                right delimiter={]},
                %below = M1,
                nodes={minimum width=3em, minimum height=3em, outer sep=0pt},
                column sep=-\pgflinewidth,
                row sep=-\pgflinewidth] (M2) { 
                                                0 & 2 & 2 & 3 & 4 & 5\\
                                                };
            \node[left] at (M2.west) {Row Index.............}; 
        \end{tikzpicture}

        \begin{tikzpicture}[>=stealth, shorten >=1pt, auto, node distance=3cm, thick, main node/.style={circle, draw, font=\sffamily\Large\bfseries}]

            \matrix[matrix of math nodes,
                left delimiter={[},
                right delimiter={]},
                %below = M2,
                nodes={minimum width=3em, minimum height=3em, outer sep=0pt},
                column sep=-\pgflinewidth,
                row sep=-\pgflinewidth] (M3) { 
                                                1 & 2 & 2 & 2 & 0 & \\
                                                };
            \node[left] at (M3.west) {Col Index...};
        
        \end{tikzpicture}
        \caption{A CSR representation.}
        \label{figure:csr}
    \end{subfigure}
    
    \caption{Four possible represenatations of the graph in Figure~\ref{figure:graph-example}}
    \label{Figure:combined-represenatations}
\end{figure}
